\uuid{z7de}
\titre{Image réciproque}
\auteur{bodin ; hanani ; mzari}
\organisation{exo7}
\licence{CC BY-NC-SA 4.0}
\datecreate{2020-01-20}
\niveau{L1}
\module{Algèbre}
\chapitre{Logique, ensemble, raisonnement}
\sousChapitre{Ensemble}
\difficulte{2}
\theme{Ensembles, applications, image réciproque}
\exo7qcmid{lille-1.ensembles.019}

\contenu{
\begin{question}
\qclassification{100.02,101.01,102.02,102.02}
\qtags{motcle=image réciproque}
On considère l'application $f:\Rr\times \Rr\to \Rr$ définie par
$$f(x,y)=x^2+y^2.$$
Quelles sont les bonnes réponses ?
\begin{answers}
    \good{$f^{-1}(\{0\})=\{(0,0)\}$}
    \bad{$f^{-1}(\{1\})=\{(1,0)\}$}
    \bad{$f^{-1}(\{0\})=\{(0,1)\}$}
    \good{$f^{-1}(\{1\})$ est le cercle de centre $(0,0)$ et de rayon $1$}
\end{answers}
\begin{explanations}
D'abord, $x^2+y^2=0\Leftrightarrow (x,y)=(0,0)$. Par ailleurs, l'ensemble des solutions $(x,y)$ de $x^2+y^2=1$ est le cercle de centre $(0,0)$ et de rayon $1$.
\end{explanations}
\end{question}
}
